\section{Kiến thức nền: tensor, phân rã và ba thành phần}
\label{sec:tensor-primer}

\subsection{Scalar, vector, ma trận, tensor}

Từ ``tensor'' nghe đáng sợ, nhưng với bài này chỉ cần hiểu là \emph{mảng nhiều chiều}. Bảng
dưới đây là toàn bộ kiến thức tối thiểu cần cho các công thức của paper.

\begin{center}
\small
\begin{tabular}{@{}llll@{}}
  \toprule
  Đối tượng & Bậc & Ký hiệu ví dụ & Cách đọc \\
  \midrule
  Scalar & 0 & $x=7$ & một số \\
  Vector & 1 & $\mathbf{w}=(w_1,\ldots,w_k)$ & danh sách theo từ khóa \\
  Ma trận & 2 & $\mathbf{W}\in\mathbb{R}^{k\times l}$ & bảng từ khóa $\times$ địa điểm \\
  Tensor & 3 & $\mathcal{X}\in\mathbb{R}^{k\times l\times L_c}$ & chồng bảng theo thời gian \\
  \bottomrule
\end{tabular}
\end{center}

Chỉ số được viết từ trái sang phải theo đúng thứ tự trục. Với tensor bậc ba, $x_{ij}(t)$ là phần
tử ở từ khóa $i$, địa điểm $j$, thời điểm $t$. Nếu $k=2,l=2,L_c=3$, tensor có $2\cdot2\cdot3=12$
phần tử. Phép ``unfold'' (trải phẳng) biến tensor thành ma trận bằng cách gộp một số trục; ví dụ
unfold theo trục từ khóa cho ta một ma trận để giải bằng đại số tuyến tính. Ta không cần tự làm
phép unfold bằng tay, nhưng cần hiểu nó chỉ là \emph{đổi cách xếp các số}, không tạo thêm dữ liệu.

\subsection{Phân rã Tucker bằng trực giác}

Giả sử ta có $k$ từ khóa và $l$ địa điểm. Thay vì lưu mọi giá trị của tensor, ta giả định dữ liệu
được tạo bởi một số ít \emph{kiểu động học} chung. Gọi $d_k$ là số nhóm từ khóa và $d_l$ là số
nhóm địa điểm, với $d_k\ll k$ và $d_l\ll l$. Khi đó, phần trend được xấp xỉ
\begin{equation}
  \widehat{\mathcal{X}}_d
  =\mathcal{W}^{(\mathrm{core})}
   \times_{\mathrm{key}}\mathbf{W}^{(\mathrm{key})}
   \times_{\mathrm{loc}}\mathbf{W}^{(\mathrm{loc})},
  \label{eq:tucker-intuition}
\end{equation}
trong đó $\mathcal{W}^{(\mathrm{core})}$ là một tensor nhỏ $d_k\times d_l\times L_c$, còn
$\mathbf{W}^{(\mathrm{key})}$ và $\mathbf{W}^{(\mathrm{loc})}$ là các bảng hệ số không âm kích
thước $d_k\times k$ và $d_l\times l$.

Ký hiệu $\times_{\mathrm{key}},\times_{\mathrm{loc}}$ gọi là \emph{mode product}. Có thể đọc
công thức~\eqref{eq:tucker-intuition} bằng lời như sau: lấy lõi nhỏ, trải nó ra theo các nhóm
từ khóa với trọng số $\mathbf{W}^{(\mathrm{key})}$, rồi trải theo các nhóm địa điểm với trọng số
$\mathbf{W}^{(\mathrm{loc})}$. Nếu một từ khóa có trọng số cao ở nhóm 2, nó sẽ nhận nhiều hơn
dynamics của nhóm 2. Nếu mọi trọng số đều không âm, việc trộn không đảo dấu của dynamics, nên
ta vẫn đọc được ``nhóm này tăng'' hay ``nhóm này giảm''.

\begin{keypoint}
Phân rã Tucker không phải là nén file tùy tiện. Nó tạo ra một lõi nhỏ chứa \emph{động học chung}
và các hệ số nói \emph{thành viên nào thuộc nhóm nào}. Đây là cầu nối giữa dự báo và diễn giải.
\end{keypoint}

\subsection{Ví dụ \texorpdfstring{$2\times2\times3$}{2x2x3} và lợi ích nén}

Hãy tưởng tượng hai từ khóa, hai địa điểm, ba tuần. Raw tensor có 12 số. Nếu giả sử chỉ có một
nhóm từ khóa và một nhóm địa điểm ($d_k=d_l=1$), lõi có $1\cdot1\cdot3=3$ số, hai hệ số có
$1\cdot2+1\cdot2=4$ số: tổng 7 tham số thay cho 12. Với dữ liệu thật, khoảng cách lớn hơn nhiều.
Một cửa sổ Device có $8\cdot50\cdot104=41{,}600$ phần tử. Với $d_k=3,d_l=2$, lõi có
$3\cdot2\cdot104=624$ phần tử và hai hệ số có $3\cdot8+2\cdot50=124$ phần tử: khoảng 748 tham số,
chưa bằng 2\% kích thước cửa sổ.

\begin{plainwords}
Nén ở đây không phải để tiết kiệm ổ đĩa. Mục đích là buộc mô hình giải thích hàng chục nghìn
quan sát bằng một vài dynamics và một bảng ``ai thuộc nhóm nào''. Nếu mô hình vẫn khớp tốt, cấu
trúc nhóm là một giả thuyết có ích.
\end{plainwords}

\subsection{Trend, seasonal và outlier trong một chuỗi nhỏ}

Ba thành phần trong D-Tracker cộng lại xấp xỉ dữ liệu:
\begin{equation}
  \mathcal{X}^{c}\approx\widehat{\mathcal{X}}_d
  +\widehat{\mathcal{X}}_s+\widehat{\mathcal{X}}_o.
  \label{eq:decompose-primer}
\end{equation}
Bảng~\ref{tab:three-components} là ví dụ số \textbf{do tài liệu tạo}. Dòng ``quan sát'' bằng
đúng tổng ba dòng còn lại; dấu gạch ngang ở outlier nghĩa là 0.

\begin{table}[htbp]
  \centering
  \small
  \caption{Ví dụ minh họa phép cộng trend + seasonal + outlier.}
  \label{tab:three-components}
  \begin{tabular}{lrrrrrrrr}
    \toprule
    Thành phần & Tuần 1 & 2 & 3 & 4 & 5 & 6 & 7 & 8 \\
    \midrule
    Quan sát      & 10 & 12 & 15 & 14 & 11 & 13 & 17 & 25 \\
    Trend         & 10 & 11 & 12 & 13 & 14 & 15 & 16 & 17 \\
    Seasonal      &  0 &  1 &  3 &  1 & $-3$ & $-2$ &  1 &  3 \\
    Outlier       &  -- & -- & -- & -- & -- & -- & -- &  5 \\
    \bottomrule
  \end{tabular}
\end{table}

Trend là phần thay đổi tương đối mượt. Seasonal là phần lặp lại có chu kỳ (ở đây minh họa bằng
một pattern ngắn; dữ liệu thật dùng chu kỳ tuần/năm). Outlier là phần bất thường và được giả định
\emph{thưa}: chỉ một số ít ô khác 0. Tách như vậy giúp mô hình không dùng một cú tăng đột biến
để kết luận rằng cả xu hướng đã thay đổi.

\subsection{Từ phân rã đến câu hỏi của paper}

Ta đã có ba mảnh: một tensor quan sát, một cách nén trend bằng lõi và hệ số, và một phép tách
thành ba thành phần. Paper thêm hai ý tưởng nữa:
\begin{itemize}
  \item lõi trend không tiến hóa tùy ý, mà tuân theo một hệ \textbf{reaction--diffusion}; và
  \item mô hình được cập nhật theo stream, với cấu trúc (số nhóm) và thời điểm đổi mô hình được
        chọn tự động.
\end{itemize}
Chương~\ref{sec:reaction-diffusion} giải thích ý tưởng thứ nhất; Chương~\ref{sec:model} ghép tất
cả lại thành mô hình hoàn chỉnh.

\begin{commonmistake}
Ba lỗi hiểu thường gặp: (1) coi $\mathcal{W}^{(\mathrm{core})}$ là tensor dữ liệu gốc thu nhỏ ---
nó là \emph{tham số ẩn} được học; (2) coi hệ số $\mathbf{W}$ là xác suất --- chúng là trọng số
không âm, không nhất thiết cộng lại bằng 1; (3) coi phép tách là duy nhất --- các phép phân rã
tensor thường chỉ xác định tới một phép biến đổi, và paper giải một bài tối ưu cụ thể.
\end{commonmistake}

\begin{selfcheck}
\begin{enumerate}
  \item Với $d_k=2,d_l=3,L_c=10$, lõi có bao nhiêu phần tử?
  \item Giải thích bằng một câu vì sao hệ số không âm giúp giữ chiều của dynamics.
  \item Trong bảng ba thành phần, nếu bỏ outlier ở tuần 8 thì mô hình sẽ hiểu sai điều gì?
\end{enumerate}
\end{selfcheck}
